\section{Marc Teòric}

\subsection{Ecosistema}
Anomenem ecosistema al sistema físic i biològic format per una comunitat d’éssers vius que habiten en un espai delimitat (per exemple un bosc). Es tracta d’un conjunt de factors físics i factors biològics. Quan parlem de factors físics ens referim a factors ``no vius'' tals com el clima, la llum solar, l’aigua, el sòl i substrat, i el relleu geològic, i quand parlem de factors biològics ens referim a tot el que considerem viu com ara els diversos microorganismes que podem trobar en el sòl, las plantes, animals\dots Els ecosistemes estan estructurats de forma que segueixen una cadena, si una de les parts es veu alterada o afectada per un nou factor com per exemple el canvi climàtic la resta de la estructura del ecosistema també es veurà afectada, i per això els ecosistema pateixen tant per el canvi climàtic, ja que són molt fràgils i al mínim canvi pot alterar greument el ecosistema o fins i tot acabar amb ell.

El concepte d’ecosistema es va començar a desarrollar entre els anys 1920 i 1930 i tenen en compte les complexes interaccions entre individus que formen una comunitat (biocenosi) i els fluxes d’energía i materials que els travesen. Aquest terme fou encunyat per Roy Clapham. Un exemple d’aquestes interaccions és la cadena tròfica.

La cadena tròfica està majoritàriament formada per 4 grups: Els productors, els consumidors primaris, els secundaris i els descomponedors. Aquesta cadena explica com els productors (plantes) produeixen energia orgànica a partir d’energia inorgànica (sol, aigua, etc.), després els consumidors primaris (animals herbívors) s’alimenten d’aquestes plantes i extreuen l’energia d’aquestes, seguidament els consumidors secundaris depredan els consumidors primaris per tal d’obtenir energia, i finalment els descomponedors que transforman i obtenen energia a partir dels residus que deixen els altres 3 grups.

\subsection{Tipus d’ecosistemes}

\subsection{Relacions intraespecífiques i interespecífiques}
Les relacions intraespecífiques i interespecífiques consisteixen en la interacció biològica (relacions entre organismes dins d’un ecosistema). La gran diferència entre aquests dos grups és que les relacions intraespecífiques es refereixen a la interacció entre individus d’una mateixa espècie (com la col·laboració en una colònia de formigues), mentre que les relacions interespecíficas es refereixen a la interacció entre individus de diferents espècies (com el cas del llop gris quan depreda un cérvol).

\subsubsection{La depredació}
La depredació és una relació de caràcter interespecífic (individus de diferents espècies), on hi ha un afectat (presa) i un beneficiat (depredador), la depredació consisteix en el traspàs d’energia en el ordre presa-depredador. Els depredadors poden trobar-se en qualsevol regne. Un mateix individu pot ser depredador d'uns éssers i presa d'uns altres, sinó hi ha cap altre espècie que pugui depredar al depredador aquella espècie se la considerara el superdepredador del ecosistema.

\subsubsection{La competència}
La competència entre diferents espècies consisteix en la interacció negativa entre dos individus o poblacions a causa d’algun limitant (territori, aliment, aigua\dots). En aquests casos una de les dues bandes serà perjudica; majoritàriament a causa d’una inferioritat en aptituds físiques com ara l’estructura corporal, capacitat d’adaptació al medi, velocitat de reproducció\dots Existeixen diversos tipos de competències però les que principalment divisare en aquest treball son:

\begin{itemize}
    \item \textbf{Competència d’explotació:} Aquesta tipus de competència s’origina quan els dos organismes lluiten per recursos limitats com ara el territori o l’aliment.
    \item \textbf{Competència aparent:} Aquesta competència s'origina per la introducció d’una tercera espècie (depredador), en aquest cas principalment les espècies depredadores lluiten entre si per tal de sobreviure (per exemple la lluita per amagatalls que els ajudin a protegir-se dels depredadors), en una escala gran com el que seria un ecosistema real podem veure com depenen de l’explotació de poblacions de preses les poblacions de depredadors també augmenten.
\end{itemize}

\subsubsection{Comensalisme}
El comensalisme consisteix en dues espècies, una s’en la beneficiada de diferents formes (protecció, nutrició\dots), mentre que a l’altre espècies no li beneficia ni perjudica aquesta relació. En el cas de la biologia no deixa de ser una relació molt semblant a la simbiosi, és crucial entendrela, ya que en els ecosistemes és una de les formes que tenen les espècies de coexsistir un amb l’altre. En aquest cas també trobem diversos tipus de comensalisme:

\begin{itemize}
    \item \textbf{Foresis:} Mobilització d’un organisme gràcies a un altre que fa de medi de transport (un bon exemple son les rèmores quan s'adhereixen als taurons per recórrer llargues distàncies).
    \item \textbf{Inquilinisme:} Quan una espècie viu dins del refugi d’una altra sense causar-li cap problema (la gamba pistola i el gobi).
\end{itemize}

\subsection{Terraris}
En aquest treball l’ús de terraris és molt important, ja que d’aquesta forma es que podem realitzar l’experiment en un ambient delimitat i controlat per mi. Hi ha diferents classificacions de terraris:

\begin{itemize}
    \item \textbf{Terraris o vivaris:} Són entorns exclosivament terrestres (serà el entorn en el que treballaré).
    \item \textbf{Paludari:} Es un sistema mixt on es barregen elements aquàtics i terrestres.
    \item \textbf{Ripari:} Aquests terraris es centren en representar les vores de rius o llacs. Es diferencia amb el aquari ja que representa zones molt poc profundes on encara poden créixer plantes terrestres.
    \item \textbf{Aquari:} Es un sistema exclusivament aquàtic, a diferència del ripari aqui no trobem cap element terrestre com plantes que creixin en zones parcialment terrestres.
\end{itemize}

\subsubsection{Terrari bioactiu}
Considerem un terrari bioactiu quan en un mateix terrari trobem una espècie primària (com per exemple un rèptil), plantes, substrat fertil, microorganismes i petits invertebrats detritívors anomenats ``equip de neteja''. La característica més marcada d’un terrari bioactiu és la presència d’un equip de neteja, aquest normalment es compon d’uns microinvertebrats anomenats col·lèmbols i altres petits insectes com els isòpodes, la seva funció és descompondre matèria en descomposició (fulles seques, i transformar-la en nutrients que les plantes puguin utilitzar per créixer. Els col·lèmbols a part de descompondre matèria en descomposició també preveuen la proliferació de fongs, i els seus excrements principalment es componen de nitrats. Els isòpodes s’ocupen de menjar la matèria més gran com ara fulles mortes, per també serveixen com a suplement alimentari per al depredador principal del terrari (sempre i quan el depredador sigui insectívor).

\subsubsection{Manteniment i factors abiòtics en el terrari}
Al tractar-se d’un terrari bioactiu, el manteniment per part meva és quasi nul, gràcies als col·lèmbols que he introduït al terrari la majoria de matèria en descomposició és eliminada i d’altre banda també empedeixen la proliferació de fongs en el terrari, i els isòpodes s’encarreguen d’alimentar-se de la matèria en descomposició més gran com ara fulles mortes, o també les femtes del gecko. Per la meva part haig de fer canvis d’aigua cada dos o tres dies i humiteja la part del terrari freda una vegada al dia. La humitat sempre està al voltant del 60 \% (menys en èpoques de muda), la temperatura la regulo amb una lampara calefactora i una manta tèrmica situada al costat calent del terrari, i en cap moment del dia baixa dels 21 graus, i per últim durant la època on faig l'experiment (estiu) deixo la lámpara d’uvb encesa de 7:30-7:45 a 19:15-19:30.

\subsection{Depredador (Dragó lleopard)}

\subsubsection{Taxonomia}
És una espècie pertanyen al regne animalia, fílum chordata, a la classe reptilia, ordre squamata, família eublepharidae, al gènere eublepharis, i té de nom científic \textit{Eublepharis macularius}.

\subsubsection{Distribució}
Els dragons lleopard salvatges habiten gairebé per tot el Pakistan, part del sud de l'Afganistan i el nord-oest de l'Índia, concretament els territoris de: Rajasthan, Panjab i Jammu i Caixmir. L'espècie també es va observar el 1904 a la regió de Khorasan a l'Iran, prop de la frontera amb l'Afganistan, però aquesta observació no ha estat reeditada i ha estat posada en dubte des de llavors. Aquests límits territorials estan poc definits, si bé sí que ho estan al nord pels contraforts d'Hindu Kush. No se sap fins a quin punt els dragons lleopards s'endinsen al desert de Thar. Els dragons lleopards habiten principalment en àrees pedregoses desèrtiques o semidesèrtiques. Es poden trobar en deserts i matollars xeròfils, estepes i boscos tropicals i subtropicals de frondoses secs. Solen preferir el sòl d'argila i evitar els deserts de sorra. És comú trobar-los en espais creats per l'home, com poden ser en murs i edificis, especialment entre les esquerdes, sota les pedres i forats a terra.

\subsubsection{Mètodes de defensa i supervivència}
Els dragons lleopards com molts altres rèptils petits, tenen diversos depredadors com ara els xacals, guineus, serps\dots Els sentits de l'oïda i la vista ens tenen molt desarrollats, cosa la qual els ajuda a sobreviure durant la nit, apart la seva coloració els ajuda a camuflarse amb l’entorn. Si se senten amenaçats emeten un xiulet d’advertència per protegir-se, també utilitzen refugis com coves o forats al terra per refugiar-se del sol, sinó que també tenen la funció de protegir al dragó dels depredadors. Com a últim recurs si un dragó és atacat per un depredador pot deixar anar la cua per tal de distreure el depredador, aquesta funció s’anomena autotomia caudal.

\subsubsection{Dieta}
\subsubsection{Comportament} 
\subsubsection{Dimorfisme} 
\subsubsection{Reproducció} 
\subsubsection{Especimen juvenil} 
\subsubsection{Especimen adult} 
\subsubsection{Malalties} 

\subsection{Presa número 1 (\textit{Porcellio laevis})}
\subsubsection{Taxonomia} 
\subsubsection{Distribució} 
\subsubsection{Característiques} 
\subsubsection{Dieta} 
\subsubsection{Comportament} 
\subsubsection{Dimorfisme} 
\subsubsection{Reproducció} 
\subsubsection{Especiment juvenil} 
\subsubsection{Especiment adult} 

\subsection{Presa número 2 (\textit{Armadillidium vulgare})}
\subsubsection{Taxonomia} 
\subsubsection{Distribució} 
\subsubsection{Característiques} 
\subsubsection{Dieta} 
\subsubsection{Comportament} 
\subsubsection{Dimorfisme} 
\subsubsection{Reproducció} 
\subsubsection{Especiment juvenil} 
\subsubsection{Especiment adult}
